\section{为什么高维球交集近似 Softmax}

现在进入第一讲的核心推导。直觉是：两个半径相同的高维球中心越远，交集体积越小；在相关参数范围内，这个衰减近似指数函数。若向量经过归一化，Hamming distance 又能改写为 cosine/dot-product similarity，于是 intersection weight 变成 softmax logit 的指数形式。

\subsection{从交集曲线到指数衰减}

章节页之后的图把横轴设为两球中心的 Hamming distance，纵轴是 expected intersection size 的对数。点列近似直线意味着原始交集大小近似指数衰减：
$$
I(\delta;d,n)\approx c_1\exp(-c_2\delta),
$$
其中 $\delta=D_H(p,q)$，$d$ 是写/读半径，$n$ 是维度，$c_1,c_2$ 由 $n,d$ 决定。这个式子不是所有距离都精确，而是在注意力实际使用的相似区域给出有效近似。

\lecturefigure{slide-23-attention-approximates-sdm.jpg}{核心问题：Transformer attention 如何近似 SDM}{官方录像恢复幻灯片 V023；00:12:26}

\lecturefigure{slide-24-exponential-circle-decay.jpg}{高维 sphere intersection 随中心距离近似指数衰减}{官方录像恢复幻灯片 V024；00:12:46}

\begin{warningbox}{二维圆图只负责直觉}
二维中两个圆的交集面积不会自动呈现高维浓缩现象。真正结论依赖 $n$ 很大、半径与维度处于合适比例。不能从画面上“看起来像指数”推出定理；课程用数值图和渐近分析共同支持近似。
\end{warningbox}

\subsection{Hamming distance 怎样变成 cosine similarity}

对 bipolar 向量 $a,b\in\{-1,+1\}^{n}$，每个相同 bit 对点积贡献 $+1$，不同 bit 贡献 $-1$，因此
$$
D_H(a,b)=\frac{n-a^{\top}b}{2}.
$$
若进一步令 $\hat a=a/\sqrt n$、$\hat b=b/\sqrt n$，则
$$
D_H(a,b)=\frac{n}{2}\left(1-\hat a^{\top}\hat b\right).
$$
把它代入指数衰减：常数项被归一化消去，剩余权重正比于 $\exp(\beta\hat a^{\top}\hat b)$。这就是从距离型联想记忆到 dot-product attention 的数学桥梁。

\lecturefigure{slide-25-hamming-cosine-derivation.jpg}{从 Hamming distance 到归一化 dot product}{官方录像恢复幻灯片 V025；00:16:24}

\lecturefigure{slide-26-exponential-fit-regression.jpg}{对 log intersection size 做线性拟合}{官方录像恢复幻灯片 V026；00:17:18}

\begin{knowledgebox}{读图：拟合图支持什么}
横轴增加时，log intersection size 近似线性下降，支持指数近似；斜率对应 effective beta，截距在 softmax 归一化中被消去。图不证明所有维度和半径都同样好，弯曲尾部正是近似误差需要检查的区域。
\end{knowledgebox}

\subsection{把两套记号逐项对齐}

Attention-in-SDM-notation 页先把 query、pattern address 与 stored value 对齐；核心 equivalence 页再把 circle-intersection weight 替换为 exponential dot product。归一化后，SDM pattern view 的读出为
$$
\hat y(q)
=\frac{\sum_\mu I(D_H(p_\mu,q);d,n)v_\mu}
{\sum_\nu I(D_H(p_\nu,q);d,n)}
\approx
\sum_\mu
\frac{\exp(\beta p_\mu^{\top}q)}
{\sum_\nu\exp(\beta p_\nu^{\top}q)}v_\mu.
$$
其中 $p_\mu$ 是第 $\mu$ 条 key/pattern，$v_\mu$ 是 value，$\beta$ 由维度、半径、向量范数和近似系数共同决定。右侧正是 attention read。

\lecturefigure{slide-27-attention-sdm-notation.jpg}{用 SDM 记号重写 attention}{官方录像恢复幻灯片 V027；00:17:20}

\lecturefigure{slide-28-sdm-attention-equivalence.jpg}{SDM intersection weighting 与 attention softmax 的对应}{官方录像恢复幻灯片 V028；00:17:22}

\subsection{数值实验与成立条件}

Parameter experiments 页比较不同 $n,d$ 下的真实 intersection 与指数拟合；Mapping conditions 页强调两个实际条件：向量需要近似 $L^2$ 归一化，常由 LayerNorm 或 norm behavior 支持；effective coefficient 要匹配 intersection slope。若 key/query norms 自由变化，它们也会改变 temperature，不能只把 $1/\sqrt{d_k}$ 当唯一 beta。

\lecturefigure{slide-29-parameter-experiments.jpg}{不同维度与半径下的 intersection 近似实验}{官方录像恢复幻灯片 V029；00:18:38}

\lecturefigure{slide-30-mapping-conditions.jpg}{映射成立所需的 normalization 与 coefficient 条件}{官方录像恢复幻灯片 V030；00:20:40}

\teachervoice{讲者没有把“约等于”藏起来。他反复区分 close approximation 与 exact identity，并指出稀疏/硬 attention、未归一化向量或不同 coefficient 都会偏离标准映射。真正可迁移的方法是先写出 assumptions，再问训练模型是否满足，而不是看见 softmax 就宣布等价。}

\begin{importantbox}{核心结论的最小版本}
在高维、合适半径、归一化向量和 coefficient 匹配的条件下，SDM 读写球的交集数量随 pattern--query 距离近似指数衰减；距离可改写为 dot product；归一化后得到 softmax attention。这个结论解释 softmax 的一种几何来源，但不是 softmax 的唯一来源。
\end{importantbox}

\subsection{本章小结}

高维 intersection concentration 产生指数距离权重，bipolar Hamming distance 又等价于归一化 dot product 的仿射变换，因此 SDM pattern read 近似 attention。数值实验支持该近似，但 LayerNorm、向量范数、beta 和参数区域决定映射质量。

